\subsection{半单交换代数的张量积}

定理 3. —— 设 \(\mathbf{Z}_1\) 与 \(\mathbf{Z}_2\) 是 \(\mathbf{K}\) 上的半单交换代数。环 \(\mathbf{Z}_1 \otimes \mathbf{Z}_2\) 的根等于它的幂零元素所成之集。

我们先处理 \(Z_{1}\) 与 \(Z_{2}\) 是体 K 的扩张 \(L_{1}\) 与 \(L_{2}\) 的情形。必要时交换 \(L_{1}\) 与 \(L_{2}\)，可化归为 \(L_{1}\) 在 K 上的超越次数小于 \(L_{2}\) 在 K 上的超越次数的情形。

选取一个代数闭包 \(\Omega\)（\(\mathrm{L}_2\) 的）；由 V, §14, No.~6, 第 114 页，定理 5 的推论 1，我们可以假设 \(\mathrm{L}_1\) 是 \(\Omega\) 的一个子扩张。

\begin{enumerate}
\def\labelenumi{\Alph{enumi})}
\item
  我们先证明 \(L_{1} \otimes L_{2}\) 的根含于 \(L_{1} \otimes \Omega\) 的根。令 \(\mathfrak{a} = \Re(L_{1} \otimes L_{2})(L_{1} \otimes \Omega)\)；它是交换环 \(L_{1} \otimes \Omega\) 的一个理想，我们必须证明 a 含于 \(L_{1} \otimes \Omega\) 的根。换句话说（VIII, 第 156 页，定理 1），我们必须证明对 \(x \in a\)，元素 \(1 + x\) 在 \(L_{1} \otimes \Omega\) 中可逆。由于 \(\Omega\) 是 \(L_{2}\) 的代数扩张，存在一个有限次数的扩张 \(L_{3}\)（\(L_{2}\) 的），使得 x 属于 \(\Re(L_{1} \otimes L_{2})(L_{1} \otimes L_{3})\)。显然只需证明 \(1 + x\) 在 \(L_{1} \otimes L_{3}\) 中可逆。而 \(C = L_{1} \otimes L_{3}\) 是环 \(B = L_{1} \otimes L_{2}\) 上的一个有限生成模。由 VIII, 第 175 页的推论，我们有 \(\Re(B)C \subset \Re(C)\)，故 x 属于 C 的根，从而 \(1 + x\) 在 C 中可逆。
\item
  我们来证明 \(L_{1} \otimes \Omega\) 的根由幂零元素组成。用 p 表示 K 的特征指数，用 P 表示 K 在 \(\Omega\) 中的相对 p-根式（即纯不可分）闭包（V, §5, No.~2, 第 25 页）；体 P 是完备的。由于 P 是 K 的代数扩张，我们有 \(\mathrm{L}_{1}(\mathrm{P}) = \mathrm{L}_{1}[\mathrm{P}]\)（V, §3, No.~2, 第 18 页，推论 1）。设 b 是从 \(L_{1} \otimes P\) 到体 \(P_{1} = L_{1}[P]\) 的典范同态的核。设 \(x \in b\)；存在元素 \(y_{1}, \ldots, y_{n}\)（\(L_{1}\) 的）与元素 \(z_{1}, \ldots, z_{n}\)（P 的），使得 \(x = \sum_{i=1}^{n} y_{i} \otimes z_{i}\) 且 \(\sum_{i=1}^{n} y_{i} z_{i} = 0\)。由于 P 在 K 上是 p-根式的，存在 p 的一个幂 q，使得 \(z_{1}^{q}, \ldots, z_{n}^{q}\) 属于 K。于是我们有
\end{enumerate}

\[
x ^ {q} = \sum_ {i = 1} ^ {n} y _ {i} ^ {q} \otimes z _ {i} ^ {q} = \sum_ {i = 1} ^ {n} y _ {i} ^ {q} z _ {i} ^ {q} \otimes 1 = \left(\sum_ {i = 1} ^ {n} y _ {i} z _ {i}\right) ^ {q} \otimes 1 = 0.
\]

所以 b 由幂零元素组成。

令 \(c = b \otimes_{P} \Omega\)；它是从 \((\mathrm{L}_{1} \otimes \mathrm{P}) \otimes_{\mathrm{P}} \Omega\) 到 \(P_{1} \otimes_{P} \Omega\) 的典范同态的核，且由上面的论证知它由幂零元素组成。而 \(\Omega\) 是 P 的代数封闭扩张，\(P_{1}\) 是 \(\Omega\) 的一个子扩张。由于体 P 是完备的，\(P_{1}\) 是 P 的可分扩张（V, §15, No.~5, 第 125 页，定理 3）。由 V, 第 120 页，定理 4，交换环 \(P_{1} \otimes_{P} \Omega\) 的诸极大理想之交约化为 0。换句话说，环 \(P_{1} \otimes_{P} \Omega\)（它同构于 \(((L_{1} \otimes P) \otimes_{P} \Omega)/c\)）是无根的。这就证明（VIII, 第 155 页，命题 5）c 包含环 \((\mathrm{L}_{1} \otimes \mathrm{P}) \otimes_{\mathrm{P}} \Omega\) 的根。而这一环同构于 \(L_{1} \otimes \Omega\)，且 c 由幂零元素组成。因此，\(L_{1} \otimes \Omega\) 的根由幂零元素组成。

\begin{enumerate}
\def\labelenumi{\Alph{enumi})}
\setcounter{enumi}{2}
\tightlist
\item
  特殊情形证明的结尾。由 A) 与 B)，根 \(\mathfrak{r}\)（\(\mathrm{L}_1\otimes \mathrm{L}_2\) 的）含于这一交换环的幂零元素所成之集 \(\mathfrak{n}\)。此外我们知道 \(\mathfrak{n}\) 含于 \(\mathfrak{r}\)（VIII, 第 157 页，注 2）。
\end{enumerate}

现在考虑一般情形。由于半单交换 K-代数是体 K 的有限多个扩张的积（VIII, 第 137 页，命题 3），且环的积的根是诸根的积（VIII, 第 156 页，推论 3），\(Z_{1} \otimes Z_{2}\) 的根是这一环的幂零元素所成之集。

